\clearpage
\subsection{二重覆盖族}\label{sec:18b}
现在构造一个 Lefschetz 双纤维化，将 $\widetilde\pi:\widetilde E\to\widetilde S$ 及其他分歧点所对应的类似纤维化包含在内，也包括退化的，即奇异的情形。令 $\mathcal W\subset\CC$ 为半径 $1-2\delta$ 的圆盘。考虑沿对角线分歧的 $\mathcal W\times S$ 的二重覆盖，以及沿 $\pi$ 的图像分歧的 $\mathcal W\times E$ 的二重覆盖；更精确地说，使用对角线或图像中相应坐标落在 $\mathcal W\subset S$ 内的部分：
\[
\begin{aligned}
\Sfamily&=\{(w,y)\in\mathcal W\times\CC:y^2+w\in S\},\\
\Efamily&=\{(w,y,x)\in\Sfamily\times E:y^2+w=\pi(x)\}.
\end{aligned}
\]
$\Sfamily$ 带有自然复结构，而 $\Efamily$ 带有殆复结构；二者的边界均弱凸。有下列伪全纯映射：
\[
\begin{aligned}
\iota_{\Sfamily}(w,y)&=(w,-y):\Sfamily\to\Sfamily,\\
\iota_{\Efamily}(w,y,x)&=(w,-y,x):\Efamily\to\Efamily
&&\text{覆盖对合；}\\
\beta_{\Sfamily}(w,y)&=y^2+w:\Sfamily\to S,\\
\beta_{\Efamily}(w,y,x)&=x:\Efamily\to E
&&\text{忘掉 $w$ 的商映射；}\\
\varpi(w,y,x)&=(w,y):\Efamily\to\Sfamily
&&\text{对每个 $w$，为 $\pi$ 的提升；}\\
\psi(w,y)&=w:\Sfamily\to\mathcal W,\\
\Psi(w,y,x)&=w:\Efamily\to\mathcal W
&&\text{向参数 $w$ 的投影。}
\end{aligned}
\]
注意 $\Psi=\psi\circ\varpi$。$\varpi$ 的奇点为
\begin{equation}\label{eq:18.6}
\begin{aligned}
\operatorname{Crit}(\varpi)&=\coprod_{x\in\operatorname{Crit}(\pi)}\{y^2+w=\pi(x)\},\\
\operatorname{Critv}(\varpi)&=\coprod_{z\in\operatorname{Critv}(\pi)}\{y^2+w=z\}.
\end{aligned}
\end{equation}
在第一个集合中，$\Psi$ 的临界点恰为形如 $(w=\pi(x),0,x)$ 的那些点，故 $\operatorname{Critv}(\Psi)=\operatorname{Critv}(\pi)$。\znote{268}{原书此处印作 $\operatorname{Critv}(\Psi)=\operatorname{Crit}(\pi)$，两侧所在的空间并不相同：临界值落在 $\mathcal W\subset\CC$ 中，而临界点是 $E$ 中的点，故应视为笔误。与之相符的正确关系是 $\operatorname{Critv}(\Psi)=\operatorname{Critv}(\pi)$，译文即按此印出，可与\eref{18.6}第二式（$\operatorname{Critv}(\varpi)=\coprod_{z\in\operatorname{Critv}(\pi)}\{y^2+w=z\}$）对照。按《译校说明》的体例，原书笔误不回改，仅在此注明。[来源：原书印刷页 268 与\eref{18.6}；原书自身前后不一致，属笔误]}若在满足 $\pi(x)=z$ 的点 $x\in\operatorname{Crit}(\pi)$ 附近选择局部全纯 Morse 坐标，则 $\Efamily$ 的定义方程成为
\begin{equation}\label{eq:18.7}
w=z-y^2+x_1^2+\cdots+x_{n+1}^2;
\end{equation}
消去变量 $w$ 后，有 $\varpi(x,y)=(z-y^2+x_1^2+\cdots+x_{n+1}^2,y)$。由此容易看出，$\varpi$ 与 $\Psi$ 的临界点都是一般的。此外，将 $\Psi$ 在临界点 $(z,0,x)$ 处的 Hessian 限制到
\SourcePage{277}
$\ker(D\varpi)$，即与 $y=0$ 相切的子空间上，仍然非退化。这些计算说明，就殆复几何而言，
\SourceEquationDiagram{diagram-18-8}{0.232615134}{18.8}
满足 Lefschetz 双纤维化的全部性质；事实上，由于没有伪临界点，这是双纤维化的一个特别简单的例子。尽管整个讨论十分初等，仍值得用更直观的方式重述要点。对每个 $w$，限制
\begin{equation}\label{eq:18.9}
\varpi_w:\Efamily_w=\Psi^{-1}(w)\longrightarrow\Sfamily_w=\psi^{-1}(w)
\end{equation}
是沿二重分歧覆盖 $\beta_{\Sfamily}|\Sfamily_w:\Sfamily_w\to S$ 对 $\pi$ 作的拉回；特别地，$w=*$ 时，就回到前面定义的 $\widetilde\pi:\widetilde E\to\widetilde S$。只要 $w\notin\operatorname{Critv}(\pi)$，$\Efamily_w$ 就光滑，且\eref{18.9}具有一般临界点；它们成对出现，对应于 $\pi$ 的各临界点。更精确地说，固定 $w$ 后，每个 $x\in\operatorname{Crit}(\pi)$ 对应于两点 $(w,\pm\sqrt{\pi(x)-w},x)$。当 $w$ 趋于临界值 $z=\pi(x)$ 时，两个平方根汇合，其汇合处就是 $\Efamily_z$ 的唯一奇点。

取使 $\mathcal W$ 成为带边界正合辛流形的 $\theta_{\mathcal W},\omega_{\mathcal W}$，并假设它们充分大。考虑如下的一形式 $\theta_{\Sfamily}$：将 $\theta_{\mathcal W}$ 沿 $\psi$ 的拉回、$\kappa$ 沿第二投影的拉回，以及 $\theta_S$ 沿 $\beta_{\Sfamily}$ 的拉回相加。令 $\omega_{\Sfamily}=d\theta_{\Sfamily}$。已知 $\omega_{\Sfamily}|\Sfamily_*=\omega_{\widetilde S}$ 为辛形式，在其他任意纤维上的限制也如此。因此，由于加入了底空间上充分大的辛形式 $\omega_{\mathcal W}$ 的拉回，$\omega_{\Sfamily}$ 本身也是辛形式。用\stref{引理}{15.3}中的相同论证，特别参见\eref{15.7}，可知负 Liouville 向量场沿位于 $\partial\mathcal W$ 上方的边界面 $\partial^v\Sfamily$ 指向内部。为处理另一边界面，考虑部分平凡化
\begin{equation}\label{eq:18.10}
\mathcal W\times\widetilde S\supset\mathcal W\times\widetilde U\longrightarrow\Sfamily,\qquad
(w,y)\longmapsto(w,\sqrt{y^2+*-w}),
\end{equation}
它定义在子集 $\widetilde U=\{|y^2+*|\ge1-\delta\}\subset\widetilde S$ 上；选择在 $w=*$ 时给出恒等映射的唯一平方根。\eref{18.10}的像是逐纤维边界 $\partial^h\Sfamily$ 的邻域。由于在 $\mathcal W\times\widetilde U$ 上有 $|\sqrt{y^2+*-w}|\ge\sqrt\delta$，该部分平凡化与给定辛结构及其原始一形式相容，具体来说，$\theta_{\Sfamily}$ 沿\eref{18.10}的拉回等于 $\theta_{\mathcal W}+(\theta_{\widetilde S}|\widetilde U)$。由此容易推出，$\Sfamily$ 是带角正合辛流形，且 $\psi$ 的水平切空间平行于 $\partial^h\Sfamily$。

相应地，令 $\theta_{\Efamily}=\varpi^*(\theta_{\mathcal W}+\kappa)+\beta_{\Efamily}^*\theta_E$。与前面相同的论证说明，若 $\omega_{\mathcal W}$ 充分大，则 $\omega_{\Efamily}=d\theta_{\Efamily}$ 除了可能在
\SourcePage{278}
$\Psi$ 的临界点附近以外都是辛形式。现在，如\eref{18.7}中那样，在局部坐标 $(x,y)$ 下考察一个临界点，则 $\omega_{\Efamily}$ 等于 $y$ 坐标中的 $d\kappa$、$x$ 坐标中的 $\omega_E$，再加上 $\omega_{\mathcal W}$ 沿 $\Psi(y,x)=z-y^2+x_1^2+\cdots+x_{n+1}^2$ 的拉回；这显然是 Kähler 形式。接下来考察 $\Efamily$ 的各个边界面。(i) 对于 $|w|=1-2\delta$ 所定义的面，即 $\Psi^{-1}(\partial\mathcal W)$，前面处理 $\partial^v\Sfamily$ 的论证适用，说明负 Liouville 向量场指向内部。(ii) 在 $|y^2+w|=1$ 所定义的面附近，即 $x\in\partial^vE$ 处，有类似于\eref{18.10}的部分平凡化。因此，Liouville 向量场的行为归约为 $\mathcal W\times\widetilde E$ 中的行为。注意，该面是 $\Efamily$ 关于 $\Psi$ 的水平边界的一部分；这一平凡化还说明，$\Psi$ 的辛联络在该面附近平凡。(iii) 剩余的面由满足 $x\in\partial^hE$ 的点 $(w,y,x)$ 组成。在这样的点处，考虑
\begin{equation}\label{eq:18.11}
D\beta_{\Efamily}:T\Efamily=T\Efamily^h\oplus T\Efamily^v\longrightarrow TE=TE^h\oplus TE^v.
\end{equation}
第二项 $T\Efamily^v$ 为沿 $\varpi$ 的纤维的切空间，因此同构地映到 $TE^v$。第一项 $T\Efamily^h$ 是它关于 $\omega_{\Efamily}$ 的正交补。它实际上也就是关于 $\beta_{\Efamily}^*\omega_E$ 的正交补，因为这两个形式之差从 $\varpi$ 的底空间拉回。因此，$D\beta_{\Efamily}$ 保持分解\eref{18.11}。由于映射 $\beta_{\Efamily}$ 与 $\partial^hE$ 横截，当前边界面上的切空间是 $T(\partial^hE)$ 在\eref{18.11}下的原像，因而包含整个 $T\Efamily^h$。这说明，与 $\varpi$ 相关的辛联络的水平切空间，在 $(w,y,x)$ 处平行于边界；显然，$\Psi$ 也因此具有相同性质。此外，$T\Efamily$ 中的负 Liouville 向量与 $TE$ 中相应向量场的差别，只是一个 $T\Efamily^h$ 分量。由于该分量与边界相切，负 Liouville 向量必指向内部。这样，已经证明 $\Efamily$ 是带角正合辛流形，$\varpi,\Psi$ 是带奇点的正合纤维化；这就验证了\eref{18.8}满足 Lefschetz 双纤维化定义中关于辛结构的要求。为便于后面引用，记录关于 $\Psi$ 的几何的一个简单观察：

\StatementNumber{1}\begin{lemma}\label{st:18.1}
设 $\lambda$ 为沿 $\mathcal W$ 中半径 $|*|=1-3\delta$ 的圆周正向绕行一周的闭路径。考虑 $\Efamily_*=\widetilde E$ 上相应的单值变换自同构。在子集 $\widetilde E^{\mathrm{shrink}}$ 上，该映射与对合 $\iota_{\widetilde E}$ 一致。
\end{lemma}
\begin{proof}
在任意满足 $|w|=1-3\delta$ 且 $|\pi(x)|\le1-4\delta$ 的点 $(w,x,y)$，有 $|y|=|\sqrt{\pi(x)-w}|\ge\sqrt\delta$，因此辛形式中 $d\kappa$ 的贡献为零。于是，$\Psi$ 的辛联络的限制，等于沿投影到 $\mathcal W\times E$ 的平凡联络的拉回。因此，沿
\SourcePage{279}
$\lambda(t)=(1-3\delta)e^{2\pi it}$ 的平行移动为
\begin{equation}\label{eq:18.12}
\begin{aligned}
h_{\lambda|[0,t]}:(*,y,x)&\longmapsto\bigl(\lambda(t),\sqrt{y^2+\lambda(0)-\lambda(t)},x\bigr)\\
&=\left(\lambda(t),y\sqrt{\frac{\pi(x)-\lambda(t)}{\pi(x)-\lambda(0)}},x\right),
\end{aligned}
\end{equation}
其中选择平方根使\eref{18.12}在 $t=0$ 时为恒等映射。当 $t$ 从 $0$ 变到 $1$ 时，$\lambda(t)-\pi(x)$ 沿原点绕行一周，因此平方根改变符号。\znote{271}{这里补两句。其一，$\lambda$ 是半径 $1-3\delta$ 的整圆，而讨论中的点满足 $|\pi(x)|\le1-4\delta$，故 $\lambda(t)-\pi(x)$ 恒不为零，绕行一周是良定义的；其二，平方根变号正是 $\widetilde S\to S$ 绕分歧点 $*$ 一周交换两叶，即对合 $\iota_{\widetilde E}$：$\widetilde E$ 上正支与负支的区分只能在单连通区域内作出。§18a 末尾关于「令 $w$ 沿半径 $1-3\delta$ 的圆周移动会使 $\widetilde L_\pm$ 互换」的提醒，是同一事实的另一面。[来源：原书 §18a、§18b（印刷页 268、271）]}
\end{proof}
